\subsection{অতিরিক্ত সংক্ষিপ্ত প্রশ্ন}
\subsubsection{Data, Information ও Database}
\begin{enumerate}
\item একই শিক্ষার্থীর Address তিনটি আলাদা খাতায় রাখা হলো।
    একটি খাতায় Address পরিবর্তন করা হলো, অন্য দুটিতে করা হলো না।
    কী সমস্যা সৃষ্টি হবে?

    \item প্রকৃত নম্বর (Marks) 85-এর পরিবর্তে ডেটাবেজে 58 লেখা হলো।
    পরে ফলাফল (Result) তৈরি করলে কী সমস্যা হতে পারে?

    \end{enumerate}

\subsubsection{Data ও Database: সমন্বিত প্রয়োগ}
\noindent\textbf{CQ ধরনের ধারণামূলক প্রশ্ন।} একটি কলেজে ভর্তি, উপস্থিতি ও ফলাফলের তথ্য তিনটি আলাদা রেজিস্টারে
রাখা হয়। একই শিক্ষার্থীর মোবাইল নম্বর তিন জায়গায় লেখা আছে।
মোবাইল নম্বর পরিবর্তনের পর শুধু ভর্তি রেজিস্টারে নতুন নম্বর লেখা হলো।

\textbf{প্রশ্ন:} এই পদ্ধতির সমস্যা কী এবং ডেটাবেজ ব্যবহার করলে
কীভাবে সমস্যা কমানো যায়?





\subsubsection{Sorting}
\begin{enumerate}
\item Query Result সাজানোর Clause কোনটি?

    \item \texttt{ORDER BY Roll;} কোন Order নির্দেশ করে?

    \item GPA বেশি থেকে কম সাজাতে কী লিখবে?

    \item Text Ascending-এর সাধারণ অর্থ কী?

    \item Date Descending-এর সাধারণ অর্থ কী?

    \item একাধিক ফিল্ডে সর্টিং করলে কোন ফিল্ডের অগ্রাধিকার বেশি?

    \end{enumerate}
